\documentclass[11pt]{article}
\usepackage{lmodern}
\usepackage[T1]{fontenc}

\usepackage{amsmath,amssymb,amsthm}
\usepackage[margin=1in]{geometry}
\usepackage[hidelinks]{hyperref}

\title{A Global Heat-Sign Inequality for the
Conductor-234446 Elliptic Curve}
\author{\small Contact: \href{mailto:admin@cosmosgeek.org}{admin@cosmosgeek.org}}
\date{Working draft}

\newtheorem{theorem}{Theorem}
\newtheorem{corollary}[theorem]{Corollary}
\newtheorem{proposition}[theorem]{Proposition}


\begin{document}
\maketitle

\begin{abstract}
For the elliptic curve $y^2+xy=x^3-x^2-79x+289$ we certify exactly two
sign changes of its weight-two newform on the positive half of the
Fricke-symmetric imaginary axis. A second-order differential operator
turns the Gaussian heat transform $H$ into a strictly positive function
with strictly positive derivatives. The sole undetermined coefficient
in the resulting inhomogeneous equation is the ordinary complex
$\Lambda_E''(1)$. The zero-count theorem and certified moment bounds
show that $H(\lambda)>0$ for $|\lambda|\ge2\cdot10^{-52}$.
We also prove positivity of a continuum of hyperbolic-secant averages,
give an exact martingale marginal test and exclude a positive dilation
law between same-conductor forms. The sign or vanishing of
$\Lambda_E''(1)$ remains open.
\end{abstract}

\section{The modular heat transform}

Let
\[
E:\quad y^2+xy=x^3-x^2-79x+289
\]
be the elliptic curve of conductor
\[
N=234446.
\]
Let
\[
f_E(z)=\sum_{n\ge1}a_n e^{2\pi inz}
\]
be its normalized weight-two newform.

Define
\[
\Lambda_E(s)=
N^{s/2}(2\pi)^{-s}\Gamma(s)L(E,s).
\]
The root number of $E$ is $+1$, so
\[
\Lambda_E(s)=\Lambda_E(2-s).
\]

Put
\[
\alpha=\frac{2\pi}{\sqrt N},
\qquad
g(t)=e^t f_E\left(\frac{ie^t}{\sqrt N}\right),
\qquad t\ge0.
\]
Equivalently,
\[
g(t)=e^t\sum_{n\ge1}a_n
             \exp(-\alpha n e^t).
\]

The completed Mellin integral, split at
$y=1/\sqrt N$ and transformed using the Fricke
functional equation, gives
\begin{equation}
\Lambda_E(1+s)
=2\int_0^\infty g(t)\cosh(st)\,dt.
\label{eq:mellin}
\end{equation}

The Fourier expansion implies that $g$ decreases
superexponentially as $t\to\infty$.
All subsequent differentiations and integrations
are therefore absolutely justified.

For $\lambda\in\mathbb C$, define
\begin{equation}
H(\lambda)=
2\int_0^\infty e^{\lambda t^2}g(t)\,dt.
\label{eq:heat}
\end{equation}
This is an entire function of $\lambda$.

Its derivatives at the origin satisfy
\begin{equation}
H^{(m)}(0)=\Lambda_E^{(2m)}(1),
\qquad m\ge0.
\label{eq:derivatives}
\end{equation}
In particular,
\[
H(0)=\Lambda_E(1),\qquad
H'(0)=\Lambda_E''(1),\qquad
H''(0)=\Lambda_E^{(4)}(1).
\]

The associated two-variable transform is
\[
\Xi_\lambda(z)=
2\int_0^\infty
e^{\lambda t^2}g(t)\cos(zt)\,dt.
\]
It satisfies the backward heat equation
\[
\partial_\lambda\Xi_\lambda
=-\partial_z^2\Xi_\lambda,
\]
with
\[
\Xi_0(z)=\Lambda_E(1+iz),
\qquad
H(\lambda)=\Xi_\lambda(0).
\]

\section{A general positive-forcing theorem}

We first prove an analytic result independent
of modular forms.

Let $g:[0,\infty)\to\mathbb R$ be continuous,
decaying faster than $e^{-At^2}$ for every
$A>0$. Assume that it has exactly two simple
zeros
\[
0<t_1<t_2
\]
and satisfies
\begin{equation}
g(t)>0\quad(0\le t<t_1),\qquad
g(t)<0\quad(t_1<t<t_2),\qquad
g(t)>0\quad(t>t_2).
\label{eq:signpattern}
\end{equation}

Put
\[
x_1=t_1^2,\qquad x_2=t_2^2,
\qquad
D=\frac{d}{d\lambda},
\]
and define $H$ by \eqref{eq:heat}.

\begin{theorem}[Global positive forcing]
\label{thm:positive}
Under \eqref{eq:signpattern}, the function
\[
F(\lambda)=(D-x_1)(D-x_2)H(\lambda)
\]
has the integral representation
\begin{equation}
F(\lambda)=
2\int_0^\infty
e^{\lambda t^2}
(t^2-x_1)(t^2-x_2)g(t)\,dt.
\label{eq:forcing}
\end{equation}
For every real $\lambda$ and every integer $m\ge0$,
\begin{equation}
F^{(m)}(\lambda)>0.
\label{eq:allpositive}
\end{equation}

In particular, $F$ is strictly log-convex
on $\mathbb R$.
\end{theorem}

\begin{proof}
Differentiation under the integral defining
$H$ gives
\[
(D-x_j)H(\lambda)
=
2\int_0^\infty
(t^2-x_j)e^{\lambda t^2}g(t)\,dt.
\]
Applying both operators proves
\eqref{eq:forcing}.

By \eqref{eq:signpattern},
\[
(t^2-x_1)(t^2-x_2)g(t)>0
\]
except at the two isolated zeros.
Since $e^{\lambda t^2}>0$ for real $\lambda$,
the integral is strictly positive.

Differentiating $m$ more times inserts the
nonnegative factor $t^{2m}$. Hence
\[
F^{(m)}(\lambda)=
2\int_0^\infty
t^{2m}e^{\lambda t^2}
(t^2-x_1)(t^2-x_2)g(t)\,dt>0.
\]

For strict log-convexity, regard
\[
d\nu(x)=
\frac{(x-x_1)(x-x_2)g(\sqrt{x})}{\sqrt{x}}\,dx
\]
as a positive measure. Then
\[
F(\lambda)=\int_0^\infty e^{\lambda x}\,d\nu(x).
\]
Cauchy--Schwarz gives
\[
FF''-(F')^2\ge0.
\]
Equality would require $x$ to be constant
$\nu$-almost everywhere, which is impossible
because $\nu$ has nontrivial interval support.
The inequality is therefore strict.
\end{proof}

\section{The surviving homogeneous term}

Assume additionally that
\begin{equation}
H(0)=0.
\label{eq:centralzero}
\end{equation}

Define
\begin{equation}
B(\lambda)=
\frac{e^{x_2\lambda}-e^{x_1\lambda}}
{x_2-x_1}.
\label{eq:B}
\end{equation}
Then
\[
B(0)=0,\qquad B'(0)=1,
\qquad
(D-x_1)(D-x_2)B=0.
\]

\begin{theorem}[Exact positive decomposition]
\label{thm:decomposition}
Under \eqref{eq:signpattern} and
\eqref{eq:centralzero}, one has, for every
real $\lambda$,
\begin{equation}
\boxed{
H(\lambda)=
B(\lambda)H'(0)+
\int_0^\lambda B(\lambda-u)F(u)\,du.
}
\label{eq:duhamel}
\end{equation}
Moreover,
\begin{equation}
\boxed{
H(\lambda)-B(\lambda)H'(0)>0
\qquad(\lambda\ne0).
}
\label{eq:strictremainder}
\end{equation}
\end{theorem}

\begin{proof}
The function $B$ is the fundamental solution
of the homogeneous differential equation
\[
(D-x_1)(D-x_2)y=0
\]
with initial data $B(0)=0$, $B'(0)=1$.

Variation of constants gives the unique
solution of
\[
(D-x_1)(D-x_2)H=F
\]
satisfying $H(0)=0$ and the prescribed
value of $H'(0)$. This is precisely
\eqref{eq:duhamel}.

For $\lambda>0$ and $0\le u<\lambda$,
\[
B(\lambda-u)>0,\qquad F(u)>0.
\]
Thus the integral is positive.

If $\lambda<0$, then $B(\lambda-u)<0$
for $\lambda<u\le0$. Reversing the orientation
of integration gives
\[
\int_0^\lambda B(\lambda-u)F(u)\,du>0.
\]
This proves \eqref{eq:strictremainder}.
\end{proof}

There is also a direct convexity proof,
which gives a second positive representation.

For $\lambda\in\mathbb R$, let
\[
A(\lambda)=
\frac{x_2e^{x_1\lambda}
      -x_1e^{x_2\lambda}}{x_2-x_1}.
\]
The affine function
\[
\ell_\lambda(x)=A(\lambda)+B(\lambda)x
\]
interpolates $e^{\lambda x}$ at $x_1,x_2$.

For $\lambda\ne0$, strict convexity of
$x\mapsto e^{\lambda x}$ gives
\[
\operatorname{sgn}
\bigl(e^{\lambda x}-\ell_\lambda(x)\bigr)
=
\operatorname{sgn}((x-x_1)(x-x_2))
\]
for $x\ne x_1,x_2$.

Since $H(0)=0$,
\begin{equation}
\boxed{
\begin{aligned}
&H(\lambda)-B(\lambda)H'(0)\\
&\quad=
2\int_0^\infty g(t)
\left[
e^{\lambda t^2}
-A(\lambda)-B(\lambda)t^2
\right]dt>0
\end{aligned}
}
\label{eq:convexrepresentation}
\end{equation}
for every $\lambda\ne0$.

Thus strict convexity and the sign geometry give the positive remainder.

\section{The complete real-zero dichotomy}

\begin{theorem}[Real zeros, counting multiplicity]
\label{thm:zeros}
Assume \eqref{eq:signpattern} and $H(0)=0$.
Then $H$ has at most two real zeros,
counting multiplicity.

Writing $c=H'(0)$, exactly one of the
following occurs:
\begin{enumerate}
\item If $c=0$, then $0$ is a zero of
exact multiplicity two and
\[
H(\lambda)>0\qquad(\lambda\ne0).
\]
\item If $c>0$, then $0$ is simple and
there is exactly one additional simple
zero in $(-\infty,0)$.
\item If $c<0$, then $0$ is simple and
there is exactly one additional simple
zero in $(0,\infty)$.
\end{enumerate}
\end{theorem}

\begin{proof}
Put
\[
Y(\lambda)=e^{-x_1\lambda}H(\lambda).
\]
Then
\[
Y'(\lambda)=e^{-x_1\lambda}(D-x_1)H(\lambda).
\]

Define
\[
Z(\lambda)=
e^{-x_2\lambda}(D-x_1)H(\lambda).
\]
By Theorem \ref{thm:positive},
\[
Z'(\lambda)=e^{-x_2\lambda}F(\lambda)>0.
\]
Thus $Z$ is strictly increasing and has
at most one real zero.

If $H$ had at least three real zeros
counting multiplicity, generalized Rolle's
theorem would give at least two zeros
of $Y'$, hence at least two zeros of $Z$.
This is impossible.

Now suppose $c=0$. Since $H(0)=0$,
\[
F(0)=H''(0)>0.
\]
Thus $0$ is exactly a double zero.
Theorem \ref{thm:decomposition} gives
$H(\lambda)>0$ everywhere else.

For the remaining cases, note that
$g(0)>0$. Gaussian localization gives
\[
H(-T)\sim g(0)\sqrt{\frac{\pi}{T}}>0
\qquad(T\to+\infty).
\]

Also $H(\lambda)>0$ for sufficiently
large positive $\lambda$. Indeed, the
negative part of $g$ is confined to
$(t_1,t_2)$, while a compact interval
strictly beyond $t_2$ contributes a
positive term growing faster than
$e^{\lambda t_2^2}$.

If $c>0$, then $H(\lambda)<0$ for
negative $\lambda$ sufficiently close
to zero. Positivity at $-\infty$
forces an additional negative zero.

If $c<0$, the analogous argument
forces an additional positive zero.

The two-zero bound proves uniqueness
and simplicity in both cases.
\end{proof}

\section{Isolation of the two sign changes}

Put $P(q)=\sum_{n\ge1}a_nq^{n-1}$, so the sign of $g(t)$ is the
sign of $P(e^{-\alpha e^t})$. The following certificate covers
$0\le q\le0.988$. The upper endpoint exceeds $e^{-\alpha}>0.944$.

\begin{proposition}\label{hyp:sign}
$P$ has precisely two roots in this interval. They are simple, one
in $(0.423,0.424)$ and one in $(0.943,0.944)$.
\end{proposition}
\begin{proof}
Use $M=3200$ and $P_M(q)=\sum_{n=1}^{M}a_nq^{n-1}$.
For $0\le q\le u=0.988$, Deligne's bound $|a_n|\le d(n)\sqrt n\le2n$
gives
\[
 |P-P_M|\le \frac{2u^M(M+1-Mu)}{(1-u)^2}.
\]
For $P'-P'_M$ use the derivative at $u$ of the positive-coefficient
majorant on the right, as a function of $u$. Both bounds are below
$10^{-5}$. The table records the minimum lower endpoint over the
accepted subintervals, rounded down.
\[
\begin{array}{cclrr}
1000q_{\rm left}&1000q_{\rm right}&\text{quantity}&
\#\text{ intervals}&\text{lower bound}\\ \hline
0&423&P_M&4&0.00077\\
423&424&-P'_M&1&4.05\\
424&943&-P_M&17&0.000046\\
943&944&P'_M&1&169.4\\
944&988&P_M&4&0.080
\end{array}
\]
The signs on either side of each derivative interval are opposite.
Monotonicity gives precisely one simple root on each.
\end{proof}

Here are integer instructions to reproduce the table. For odd primes
$p\le M$ use $a_p=-\sum_{x\bmod p}(\frac{4x^3-3x^2-316x+1156}{p})$;
$a_2=-1$. If $n=pm$ is composite, take the least prime factor $p$ and
$a_n=a_pa_m-p\,a_{m/p}$ when $p\mid m$, $p\nmid N$, and $a_n=a_pa_m$
otherwise. Start at $a_1=1$. For the general Weierstrass coefficients
$[a_1',a_2',a_3',a_4',a_6']$ used below the quadratic-character argument
is $(a_1'x+a_3')^2+4(x^3+a_2'x^2+a_4'x+a_6')$; at 2 use $2$ minus the
number of affine solutions modulo 2.

For the polynomial interval step use the following pseudocode.
All arrays contain coefficients in ascending order. Variables $L,U,l,h$
represent endpoints with common denominator $D$.
\begin{verbatim}
Horner(c,l,h,D):
    L,U = 0,0
    for z in reverse(c):
        L = z*D + floor(L*(l if L>=0 else h)/D)
        U = z*D + ceil (U*(h if U>=0 else l)/D)
    return [L,U]
\end{verbatim}
Use $D=10^{22}$. For each input interval apply Horner to the signed
polynomial indicated in the table (equivalently negate the returned
interval as necessary). Accept when the lower bound exceeds $10^{-5}$.
Otherwise subdivide at $\lfloor(l+h)/2\rfloor$. No floating-point
rounding enters this certificate.

Because
\[
q(t)=e^{-\alpha e^t}
\]
decreases strictly from $q_0$ to $0$,
Proposition~\ref{hyp:sign} gives two
simple zeros
\[
t_1=
\log\left(-\frac{\sqrt N}{2\pi}\log r_+\right),
\qquad
t_2=
\log\left(-\frac{\sqrt N}{2\pi}\log r_-\right).
\]
They satisfy $0<t_1<t_2$.

Since $P(0)=1$, their sign pattern is
exactly
\[
g(t):\qquad +,\ -,\ +.
\]

The exact central value is
\begin{equation}
L(E,1)=0.
\label{eq:Ecentral}
\end{equation}
This follows from the four independent rational generators recorded for this curve and the analytic-rank-zero theorem stated by Wiles \cite{lmfdb,wiles}.
It implies
\[
\Lambda_E(1)=\Lambda_E'(1)=0
\]
by the functional equation.

\begin{corollary}[Curve-specific global inequality]
\label{cor:curve}
Put
$x_j=t_j^2$. Then, for every real $\lambda$,
\begin{equation}
\boxed{
(\partial_\lambda-x_1)
(\partial_\lambda-x_2)H(\lambda)>0.
}
\label{eq:curvepositive}
\end{equation}
Every derivative of the left-hand side
is also strictly positive.

Moreover,
\begin{equation}
\boxed{
H(\lambda)=
\Lambda_E''(1)
\frac{e^{x_2\lambda}-e^{x_1\lambda}}
{x_2-x_1}
+R_E(\lambda),
}
\label{eq:curveidentity}
\end{equation}
where
\[
R_E(\lambda)=
\int_0^\lambda
\frac{e^{x_2(\lambda-u)}
      -e^{x_1(\lambda-u)}}
{x_2-x_1}
F(u)\,du
\]
satisfies
\[
R_E(\lambda)>0\quad(\lambda\ne0),
\qquad
R_E(0)=R_E'(0)=0.
\]

If $\Lambda_E''(1)=0$, then
\[
H(\lambda)>0\quad(\lambda\ne0),
\qquad
\Lambda_E^{(4)}(1)>0.
\]
In particular, the original complex
$L$-function then has a zero of
exact order four at $s=1$.
\end{corollary}

\begin{proof}
Proposition~\ref{hyp:sign} gives
\eqref{eq:signpattern}. Equation
\eqref{eq:Ecentral} gives $H(0)=0$.
Theorems \ref{thm:positive},
\ref{thm:decomposition}, and
\ref{thm:zeros} apply directly.

If $\Lambda_E''(1)=0$, then
\[
\Lambda_E^{(4)}(1)=F(0)>0.
\]
Since the gamma and conductor factors
in $\Lambda_E$ are analytic and nonzero
at $s=1$, the complex $L$-function
has exact order four.
\end{proof}

\section{A global bound on possible negative heat}

Use the even Fricke extension of $g$ to $\mathbb R$. Write
$\delta=H'(0)$ and $M_j=\int_{\mathbb R} t^jg(t)\,dt$.
In particular $\delta=M_2$ and $H''(0)=M_4$. We obtain
\begin{equation}\label{eq:micromom}
|\delta|<8\cdot10^{-49},\qquad 16541<M_4<16542.
\end{equation}
We give the details because a small value and an exact zero have
different implications for the rank.

Take $h=1/25$, $J=250$, and $M=10000$. For $j=2,4$ the finite sums
\[
D_j=2h\sum_{i=1}^J (ih)^j e^{ih}
        \sum_{n=1}^M a_n\exp(-\alpha n e^{ih})
\]
lie in $(-1.216,-1.215)10^{-60}$ and $(16541.5867,16541.5869)$,
respectively. Here $i$ is an integer index. For the infinite trapezoid
rule use $\eta=\pi/4$. On the lines $t\pm i\eta$,
\[
 |g(t\pm i\eta)|\le 2e^{-|t|}/(\alpha\cos\eta)^2.
\]
Indeed $|a_n|\le2n$ and $\sum n e^{-ny}=(2\sinh(y/2))^{-2}\le y^{-2}$;
Fricke symmetry covers $t<0$. The strip bound for the trapezoidal
rule follows by shifting the Fourier integral to the two boundary lines
and summing its nonzero frequencies \cite{TW}.
For $j=2$ the error is bounded by
\[
\frac{8(2+\eta^2)}{\alpha^2\cos^2\eta\,[\exp(2\pi\eta/h)-1]}
<6.56\cdot10^{-49}.
\]
For $j=4$ replace $2+\eta^2$ by $24+4\eta^2+\eta^4$.
At each positive mesh node, with $q=e^{-\alpha e^{ih}}$,
the omitted Fourier contribution is bounded by
\[
4h(ih)^j e^{ih}q^{M+1}
       \frac{M+1-Mq}{(1-q)^2}.
\]
The sums of these bounds are below $5\cdot10^{-57}$ and
$8\cdot10^{-60}$ respectively. Beyond $ih=10$ use
$|g(t)|\le 2e^t e^{-\alpha e^t}/(1-e^{-\alpha e^t})^2$.
For all moments through degree 24 the sum over the omitted nodes
is below $10^{-80}$ (bound consecutive terms using
$\exp(-(e^h-1)\alpha e^{10})$). These inequalities prove
\eqref{eq:micromom}.

For reproducibility, $D_j$ can be evaluated by the integer Horner
algorithm above with $D=10^{90}$: at each node enclose $e^{ih}$ and
$q$ between neighboring grid multiples, apply Horner to
$(a_1,\ldots,a_M)$, multiply outwards by $q\,e^{ih}$, then sum using
the exact rational weight $2h(ih)^j$. In forming the exponential
arguments and tail bounds use at least 100 significant decimal digits
with outward rounding. For $\pi$ one may use
$16\arctan(1/5)-4\arctan(1/239)$ and the alternating series,
and bound $\sqrt N$ by squaring rational endpoints.

\begin{corollary}
For real $|\lambda|\ge2\cdot10^{-52}$ one has $H(\lambda)>0$.
For every real $\lambda$ one has $H(\lambda)>-2.1\cdot10^{-101}$.
\end{corollary}
\begin{proof}
For $|\lambda|\le \alpha e^2/8$, Taylor's theorem gives
\[
 |H(\lambda)-\lambda\delta-\tfrac12 M_4\lambda^2|
 \le 1920\alpha^{-2}|\lambda|^3.
\]
Indeed $\exp(|\lambda|t^2)\le\exp(\alpha e^t/2)$ for $t\ge0$;
$|g(t)|\exp(\alpha e^t/2)\le8e^{-t}/\alpha^2$ by the coefficient
majorant. Integrating $t^6$ against this bound gives the stated
constant. At $\lambda=\pm 2\cdot10^{-52}$ \eqref{eq:micromom}
makes the quadratic term exceed the linear and error terms.
Theorem~\ref{thm:zeros} then implies positivity outside this interval:
in either simple-zero case the negative values form a single interval
adjacent to zero. Inside, completing the square bounds $H$ below by
$-\delta^2/(2M_4)$ minus the cubic error, giving the second claim.
\end{proof}

\section{Positive test measures and their limits}

Write $d\mu(x)=g(\sqrt x)dx/\sqrt x$, with $x=t^2$, and denote its
two sign-change points by $a=x_1$, $b=x_2$. For $K\ge0$ put
\[
 P_0(K)=\int(K-x)_+\,d\mu(x),\qquad C_0(K)=\int(x-K)_+\,d\mu(x).
\]
Then $C_0-P_0=\delta$. The first function is positive on $(0,a]$,
concave on $[a,b]$, and $C_0(K)>0$ for $K\ge b$. Also $C_0(K)\to0$.
Consequently $P_0(K)<0$ for some $K$ exactly when $\delta>0$.
Taylor's integral formula and Laplace transformation give
\[
 H(\lambda)=\lambda\delta+
       \lambda^2\int_0^\infty e^{\lambda K}C_0(K)\,dK,\qquad
 H(-u)=u^2\int_0^\infty e^{-uK}P_0(K)\,dK\quad(u>0).
\]
In particular $H(-u)/u^2$ is completely monotone if $\delta=0$.
For this curve $C_0(64)>9\cdot10^{-16}$ by the interval calculation
below. It follows from concavity and the monotonicity of $C_0$ beyond
$b$ that every $K\in(0,64]$ has $P_0(K)>0$.

A nonnegative exponential mixture $\phi(x)=\int e^{\lambda x}d\nu(\lambda)$
gives a valid test: whenever the integrals are absolutely convergent,
$\int\phi\,d\mu=\int H\,d\nu$. Nonconstant polynomials cannot be
such mixtures on $[0,\infty)$: at an interior point the even derivative
of order above the polynomial degree forces $\nu$ to be supported at zero.
The rational test $\phi(x)=x/(x+3)^2$ has an inverse Laplace density
$(1-3u)e^{-3u}$ of both signs. Cutting this density off at $u=1/3$
produces an admissible test
$\phi_+(x)=(x+3e^{-(x+3)/3})/(x+3)^2$.

Here are outward enclosures on the actual curve, with the original heat
parameter throughout:
\[
\begin{array}{ll}
-4.40<\int \phi\,d\mu<-4.35,&2.84<\int\phi_+\,d\mu<2.91,\\
2.27<\int (10+x)^{-1}d\mu<2.32,&
50.48<\int (1+x)^{-1}d\mu<50.65,\\
43.57<\int e^{-\sqrt x}d\mu<43.73,&-39.32<\Lambda_E(1+2i)<-38.77,\\
55.68<H(-1/4)<55.88,&49844<H(1/4)<50061,\\
-70.35<H(0)-2H(-1)+H(-2)<-70.00.&
\end{array}
\]
For example $(1+x)^{-1}=\int_0^\infty e^{-u}e^{-ux}\,du$ and
$e^{-\sqrt x}=\int_0^\infty e^{-ux-u^{-1}/4}\,du/(2\sqrt\pi u^{3/2})$.
The discrete difference uses weights of both signs. The negative
critical-line value is also insufficient: $\mu=\delta_0-2\delta_1+\delta_2$
has $H(\lambda)=(e^\lambda-1)^2$ while its cosine transform at frequency
$2\pi$ equals $\cos(2\pi\sqrt2)-1<0$. Gaussian Fourier inversion
specifically gives the centered averages
\[
H(-u)=(4\pi u)^{-1/2}\int_{\mathbb R}e^{-y^2/(4u)}\Lambda_E(1+iy)\,dy.
\]

\paragraph{Interval integrations.}
Split $[0,9]$ into intervals of length $1/2000$.
On each $[l,h]$ enclose $q$ between $\exp(-\alpha e^h)$ and
$\exp(-\alpha e^l)$; enclose $g$ by interval products
$[e^l,e^h]\,[q]\,P_{3200}([q])$ with the polynomial truncation
bound in Proposition~\ref{hyp:sign} evaluated at the local upper
endpoint. Multiply by the interval extension of each displayed
weight and sum twice the interval lengths times these products.
Use outward grid $10^{-38}$ in Horner and at least 100 decimal
digits for other elementary bounds. A shorter coefficient sum of
length $L=\min(3200,\max(12,\lfloor72/(\alpha e^l)\rfloor+2))$
gives the same enclosures if the corresponding tail is added at each
step. For the cosine use the midpoint value with error $1/2000$
by Lipschitz continuity in its argument $2t$.
For $C_0(64)$ integrate $(t^2-64)g(t)$ on $[8,9]$; the remaining
tail is positive. For the absolute first moment and a uniform bound
the same grid gives
\[
-108.06<\int_{\mathbb R}|t|g(t)\,dt<-106.79,\qquad \sup|g|<45.01.
\]
The absolute tails of the integrals beyond 9 are below $10^{-23}$:
there $t^2/4\le\alpha e^t/3$ and $1+t^2\le\alpha e^t$; use the
Fourier majorant above, allowing also the factor $e^{t^2/4}$.
This also bounds $\sup_{t>9}|g(t)|$.

\paragraph{Comparison at conductor 4069.}
For the curve $y^2+y=x^3-4x$ of conductor $4069$, the same integral
procedure, now with $h=1/500$ and $M=2000$, gives
$-3.46<H(-1/4)<-3.37$. Thus a negative heat value can occur for an actual
elliptic newform. For the rank-four target the positive central-sign
information is much stronger than an analogy of root numbers.

\paragraph{A family of Gamma averages.}
For $p,k>0$ set $I_p(k)=\int_{\mathbb R}\mathrm{sech}^p(kt)g(t)\,dt$.
The product for $\cosh$ \cite{DLMF} identifies
$\mathrm{sech}^p(k\sqrt x)$ as $\mathbb E e^{-Sx}$, with $S$ the
sum of independent Gamma variables of shape $p$ and scales
$k^2/(\pi^2(j+1/2)^2)$, $j\ge0$. Their means sum to $pk^2/2$.
Consequently $I_p(k)=\mathbb E H(-S)$. If $\delta\le0$, strict convexity
and the negative secant slope of the kernel give $I_p(k)>0$ for every
$p,k$. Conversely $I_p(k)=-(p/2)\delta k^2+O_p(k^4)$ as $k\downarrow0$.
Thus a negative member exists exactly when $\delta>0$.

For this curve a negative witness in that family would require $k<2/3$.
Indeed let $K_k(v)=\int e^{-v(\cosh(kt)-1)}g(t)\,dt$ for $v\ge0$.
It begins at zero with right derivative $-\Lambda_E(1+k)<0$ when
$k\ge2/3$, by the absolutely convergent elliptic Euler product, and
is positive for sufficiently large $v$ by localization at $t=0$.
The pushforward measure under $\cosh(kt)-1$ retains the two sign changes.
Three vanishing Laplace moments at distinct $v\ge0$ are impossible:
an element of the span of the three exponentials vanishing at the two
sign-change points has the sign profile of the measure, up to a global
sign. The two zeros are simple by the two-zero bound for three
exponentials, obtained by Rolle's theorem after division by the first
exponential. Thus $K_k$ has a single positive zero, changing from
negative to positive there. Next,
\[
 \Gamma(p)I_p(k)=\int_0^\infty v^{p-1}e^{-v}K_k(v)\,dv.
\]
Increasing $p$ from $p_0$ multiplies the kernel by the increasing factor
$v^{p-p_0}$, so positivity at $p_0$ implies positivity at every $p>p_0$.
For small positive $p$, $I_p(k)=-p\int\log\cosh(kt)\,g(t)\,dt+O_k(p^2)$.
The integral here is at most
\[
 -106.79k+\frac{45.01\pi^2}{12k}<0\qquad(k\ge2/3),
\]
by $\log\cosh y=|y|-\log 2+\log(1+e^{-2|y|})$ and
$\int_0^\infty\log(1+e^{-v})\,dv=\pi^2/12$.
This proves $I_p(k)>0$ for all such $k$ and every $p>0$.

\paragraph{Two-branch transport.}
Write $\mu=L\nu_L-M\nu_M+R\nu_R$ as positive lobe masses times
probabilities, with $L+R=M$. Let $\ell,m,r$ be the three means.
For $a<x<b$ the probability kernel
\[
 K(x,dy)=\frac{r-x}{r-\ell}\nu_L(dy)
          +\frac{x-\ell}{r-\ell}\nu_R(dy)
\]
has mean $x$. Its marginal after integration against $M\nu_M$
is $(L+\delta/(r-\ell))\nu_L+(R-\delta/(r-\ell))\nu_R$,
since $\delta=L\ell+Rr-Mm$. Thus a martingale with the specified
two marginals exists exactly when $\delta=0$: this kernel supplies
sufficiency, and equality of means supplies necessity. For real
$\lambda$, Jensen's inequality gives the exact decomposition
\[
 H(\lambda)=J(\lambda)+\frac{\delta}{r-\ell}
   \left[\int e^{\lambda y}\,d\nu_R-\int e^{\lambda y}\,d\nu_L\right],
 \qquad J(\lambda)=M\int\!\left[\int e^{\lambda y}K(x,dy)-e^{\lambda x}\right]d\nu_M\ge0 .
\]
More restrictive axis shifts have a concrete limitation. A symmetric
shift from $t$ to $t\pm d$ with the mean of the squares fixed must give
$t-d$ weight $1/2+d/(4t)$. If $|t-d|<t_1$ on $(t_1,t_2)$, then
$t+d<2t_2+t_1$, so the rightmost unbounded lobe cannot be its marginal.
Similarly any fixed nonnegative weights on maps $t\mapsto\pm t+d_j$
preserving the squared barycenter for an interval must have
$\sum_j w_jd_j^2=0$ by comparing constant terms.

\section{Dilation and Hecke tests}

The two Fricke-odd comparison curves of the same conductor used here have
Weierstrass coefficients
\[
 A=[1,1,0,-696,6784],\qquad C=[1,1,1,-949,-7845].
\]
Use $g_A,g_C$ with the same $\alpha$ as $g_E$ and put
$J_\theta=(1-\theta)g_C+\theta g_A$, $k_\theta=-J_\theta'$.
Suppose a positive finite scale law gives
\[
 g_E(t)=\int_{(0,\infty)} k_\theta(t/a)\,d\sigma(a).
\]
We require absolute integrability for any moment interchange below.
For $n\ge2$ write
\[
 s_n(\theta)=
 \frac{\int_{\mathbb R} t^{2n}g_E(t)\,dt}
 {2n\int_{\mathbb R}t^{2n-1}J_\theta(t)\,dt}
\]
when the denominator is nonzero. Integration by parts and $t=au$ give
$s_n(\theta)=\int a^{2n+1}d\sigma$. Thus for every polynomial $p(X)=\sum c_i X^i$
one must have
\[
 \sum_{i,j}c_ic_j s_{2+i+j}(\theta)
  =\int a^5p(a^2)^2\,d\sigma(a)\ge0 .
\]
For $\theta=1/2$ and $(c_0,\ldots,c_5)=(-106,2507,-14415,31506,-29488,10000)$,
direct interval evaluation gives the left side in $(-3.680,-3.678)$.
This excludes the half-mixture, including laws with continuous or singular
support satisfying the moment bounds.

To reproduce the calculation use the coefficient algorithm in
Proposition~\ref{hyp:sign} for each of the three curves, with $M=6500$,
$h=1/40$, $J=400$, and integer interval denominator $10^{60}$.
Calculate the moments of orders $2n$ and $2n-1$, $2\le n\le12$,
by the same doubled half-axis quadrature as in \eqref{eq:micromom};
the parity of $t^{2n-1}J_\theta(t)$ is even. Include the absolute
Fourier tail at every node before forming the intervals. In the strip
error bound use $\eta=0.8$ and the numerator
$8e^\eta j!/(\alpha^2\cos^2\eta)$ for moment order $j$, since
$\int_0^\infty(t+\eta)^j e^{-t}dt\le e^\eta j!$.
The strip error and omitted nodes together are below $10^{-39}$ at each
of these orders. Interval division gives
$0.61482<s_2(1/2)<0.61484$ and the stated separating quadratic form.
The Gram test uses moments starting with order four, so no sign of
the unknown second derivative enters.

For $\theta=3/4$ a different obstruction survives. Fourier evaluation gives
\[
 -60<g_E''(0)<-59,\qquad
 -86<g_A'''(0)<-85,\qquad48<g_C'''(0)<50.
\]
For $g_E''(0)$ evaluate
$\sum a_ne^{-\alpha n}(1-3\alpha n+(\alpha n)^2)$;
for order three use $1-7z+6z^2-z^3$, $z=\alpha n$.
The same $M=6500$ with tail bounded by the sum of the absolute monomials
times $2n e^{-\alpha n}$ suffices. Thus $k_{3/4}''(0)>0$. Boundedness of
$k_{3/4}''$ would permit differentiation under any scale law with
$\int a^{-2}d\sigma<\infty$, contradicting the sign of $g_E''(0)$.

The inverse-moment requirement is material. For any even real-analytic
rapidly decreasing $k$ with $k(0)>0$ and $k''(0)>0$, extend its Mellin
transform from positive arguments to $-2<z<0$ by
\[
 K(z)=\frac{k(0)}{z}+\int_0^1t^{z-1}(k(t)-k(0))\,dt
                       +\int_1^\infty t^{z-1}k(t)\,dt .
\]
The limits at $-2$ from the right and at 0 from the left are $+\infty$
and $-\infty$. Pick a zero $-\beta$, $0<\beta<2$. Substitution in
$\int_0^1 k(t/a)a^{\beta-1}da$ gives, near zero,
$k(0)/\beta-k''(0)t^2/(2(2-\beta))+O(t^4)$.
The full even power series converges there by termwise integration of
the analytic germ after the Mellin constant is canceled. This positive
scale law reverses the sign of the curvature. Existence for the actual
three-quarter mixture with infinite inverse-square moment remains open.

For scales in $(0,1]$ there is an exact completion criterion. Put
$r_m=s_{m+2}$. The Hausdorff conditions $(-1)^j\Delta^j r_m\ge0$
for all $j,m\ge0$ give a unique representing measure $\nu$ on $[0,1]$.
To reconstruct a finite $\sigma$ one additionally requires
$\nu(\{0\})=0$ and $\int x^{-5/2}d\nu(x)<\infty$, since $\nu$ is the
pushforward of $a^5d\sigma$ under $x=a^2$. Whenever the indicated
moments exist and all ratios are defined this reconstructs their
scale representation. In this modular setting both comparator first
central derivatives vanish (rank three and the rank-one theorem
\cite{lmfdbcomp,wiles}). The corresponding two entire heat transforms
then agree from degree two onwards and vanish at zero. Their affine
difference tends to zero on the negative half-axis by dominated
convergence, since bounded scales admit a common integrable
majorant using rapid decay. Thus their first derivatives must agree.
Equivalently, the ratio of the Mellin transforms of $g_E$ and $k_\theta$
for real $z>0$ must extend to a completely monotone function with finite
right limit at zero. Writing $a=e^{-u}$ gives necessity and the
Laplace representation theorem and Mellin uniqueness give sufficiency
for a finite measure supported on $(0,1]$.

A nonnegative \emph{pointwise} multiplier $g_E=w(t)k_\theta(t)$
with finite $w$ would require $k_\theta$ to change sign at both target
zeros. For the general linear combination $g_C+\rho g_A$ the two
required values of $\rho$ lie in $(-0.209,-0.206)$ and
$(-0.382,-0.379)$, respectively. To check this, bracket the roots
in $q$ by $(0.9430879943,0.9430879945)$ and
$(0.4239562330,0.4239562333)$ using the same signed endpoint test;
evaluate $g_X'(t)/(e^tq)=\sum a_nq^{n-1}+
\log(q)\sum n a_nq^{n-1}$ through $n=1200$ with the corresponding
geometric derivative tails. The intervals are disjoint and $g_A'$
itself is nonzero at both roots. A different proposed positive transfer
also fails immediately: if $H(0)=0$ is a positive mixture of
$K_j(\lambda-b)$, where every $K_j$ is positive off zero and vanishes
at zero, all the shift measures must be supported at $b=0$.

Finally let
\[
 J_d(\lambda)=\int_{\mathbb R} e^{\lambda t^2}g(t+d)\,dt .
\]
For $d=\log13$ or $\log35$ one has $t_1<d<t_2$. Gaussian localization
then makes $J_d(-u)<0<H(-u)$ for large $u$. The two signs already occur in a single Hecke orbit:
$a_{13}=-6$ while $a_{35}=a_5a_7=20$. The diagonal terms of these
Hecke operators give the displayed translated integrals; the other
terms must be retained before multiplying $H$ by the eigenvalue.
For every $d$, translation and parity give exactly
\[
 J_d'(0)=H'(0),\qquad
 J_d''(0)=H''(0)+6d^2H'(0).
\]
The shifted fourth moment therefore determines the missing coefficient
only if an independent signed comparison is available.

A similar centering issue occurs on removing the completed factors.
With $\ell=\log(1/\alpha)-\gamma$ one has
\[
 \Lambda_E^{(4)}(1)=\alpha^{-1}
 \{L_E^{(4)}(1)+(\pi^2-6\ell^2)L_E''(1)\}.
\]
This follows by expanding the gamma and conductor factors and using
the even functional equation and the two known central vanishings.
For an even-sign twist with nonzero central value its analogous
second derivative contains
$\alpha_D^{-1}(\pi^2/6-\ell_D^2)L(E^D,1)$ in addition to
$\alpha_D^{-1}L''(E^D,1)$, with $\alpha_D=2\pi/\sqrt{N_D}$ and
$\ell_D=\log(1/\alpha_D)-\gamma$. Its contribution to the linear
heat term need not be small.

\section{Scope of the result}

The decomposition
\[
H(\lambda)=B(\lambda)\Lambda_E''(1)
           +R_E(\lambda)
\]
is exact, and the remainder is strictly
positive at every nonzero real parameter.

Nevertheless, its positive forcing term
does not determine the homogeneous
coefficient $\Lambda_E''(1)$.

Indeed, the general theorem permits
all three possibilities
\[
\Lambda_E''(1)<0,\qquad
\Lambda_E''(1)=0,\qquad
\Lambda_E''(1)>0.
\]
In the two nonzero cases, the heat transform
has exactly one further real zero, on the
side predicted by Theorem \ref{thm:zeros}.

Therefore the sign certificate and the
global positive-forcing representation,
by themselves, do not establish the
missing second complex zero.

The sign certificate, zero classification and finite arithmetic enclosures
above localize any possible negative heat interval. The positive-measure
and scale tests specify several mechanisms within which this question can
be decided or excluded.
An additional arithmetic or automorphic
identity eliminating the surviving
homogeneous term would be required to
deduce analytic rank four for the curve.

\begin{thebibliography}{9}
\bibitem{lmfdb} The LMFDB Collaboration, \emph{Elliptic curve 234446.a1}, \href{https://www.lmfdb.org/EllipticCurve/Q/234446/a/1}{Curve record}.
\bibitem{wiles} A. Wiles, \emph{The Birch and Swinnerton-Dyer conjecture}, Clay Mathematics Institute, \href{https://www.claymath.org/wp-content/uploads/2022/05/birchswin.pdf}{Official problem description}.
\bibitem{lmfdbcomp} The LMFDB Collaboration, \emph{Elliptic curves of conductor 234446}, \href{https://www.lmfdb.org/EllipticCurve/Q/234446/}{Conductor table}.
\bibitem{TW} L. N. Trefethen and J. A. C. Weideman, \emph{The exponentially convergent trapezoidal rule}, SIAM Review 56 (2014), 385--458. \href{https://doi.org/10.1137/130932132}{doi:10.1137/130932132}.
\bibitem{DLMF} NIST Digital Library of Mathematical Functions, \emph{Infinite products and partial fractions}, \url{https://dlmf.nist.gov/4.36}.
\end{thebibliography}
\end{document}
